\documentclass[12pt]{article}
\usepackage{mc}
\mcheader{Grades 4--5}

% writing space: \space{height in cm}
\newcommand{\workbox}[1]{%
  \begin{tikzpicture}\draw[box] (0,0) rectangle (17.6,-#1);\end{tikzpicture}%
}

% pair grid: \pairgrid{max}{cell}
\newcommand{\pairgrid}[2]{%
  \begin{tikzpicture}[baseline=(current bounding box.north)]
    \foreach \i in {0,...,#1}{%
      \fill[black!8] ({(\i+1)*#2},0) rectangle ++(#2,-#2);
      \draw[cell] ({(\i+1)*#2},0) rectangle ++(#2,-#2);
      \node at ({(\i+1.5)*#2},{-0.5*#2}) {\i};
      \fill[black!8] (0,{-(\i+1)*#2}) rectangle ++(#2,-#2);
      \draw[cell] (0,{-(\i+1)*#2}) rectangle ++(#2,-#2);
      \node at ({0.5*#2},{-(\i+1.5)*#2}) {\i};
      \foreach \j in {0,...,#1}{%
        \draw[cell] ({(\j+1)*#2},{-(\i+1)*#2}) rectangle ++(#2,-#2);
      }
    }
    \node[anchor=south] at ({(#1+3)*#2/2},0.05) {other pile};
    \node[anchor=south, rotate=90] at (-0.05,{-(#1+3)*#2/2}) {one pile};
  \end{tikzpicture}%
}

\begin{document}
\fontsize{12}{16.5}\selectfont

Each game is for two players and a pile of counters. The players take turns taking counters from the pile, and each problem gives the rule for how many counters a turn may take. Whoever takes the last counter wins.

\vspace{10pt}
\problem Each turn takes 1 or 2 counters. For every starting pile from 1 to 15 counters, decide whether you would rather go first or second, and write 1st or 2nd under the number. Then decide what you would choose with 100 counters, and explain why your way of playing wins whatever your partner does.

\vspace{4pt}
\numtable{1}{15}{15}{1.17}{0.7}{1.0}{0.4}

\vspace{6pt}
\workbox{4.1}

\vspace{14pt}
\problem A pile is called a winning pile if the player about to move can be sure to win from it, and a losing pile if not. Now each turn takes 1, 3, or 4 counters. Write W or L under each pile from 1 to 30. Describe the pattern, and explain why it continues forever.

\vspace{4pt}
\numtable{1}{30}{15}{1.17}{0.7}{1.0}{0.35}

\vspace{6pt}
\workbox{3.9}

\newpage
\problem Here are seven rules. For each rule, circle the losing piles from 1 to 24. Some of the rules have exactly the same losing piles. Find them, and for one such pair of rules explain why their losing piles are the same.

\vspace{6pt}
\begin{tikzpicture}
  \foreach \t [count=\r from 0] in {{1 or 2},{1 or 3},{1, 2, or 3},{1 or 4},{1, 3, or 5},{1, 2, or 4},{1, 4, or 6}}{%
    \node[anchor=west] at (0,{-0.4-\r*1.15}) {\t};
    \foreach \k in {1,...,24}{%
      \draw[cell] ({2.6+(\k-1)*0.625},{-\r*1.15}) rectangle ++(0.625,-0.8);
      \node at ({2.6+(\k-0.5)*0.625},{-0.4-\r*1.15}) {\k};
    }
  }
\end{tikzpicture}

\vspace{8pt}
\workbox{4.3}

\vspace{14pt}
\problem Kai always takes as many counters as the rule allows. With the rule 1, 3, or 4, list every winning pile from 1 to 30 where Kai's move leaves his partner a winning pile. For which of the seven rules in Problem 3 does Kai's move from a winning pile always leave a losing pile?

\vspace{6pt}
\workbox{4.6}

\newpage
\problem A rule is a list of how many counters a turn may take. For each list of losing piles below, find a rule that includes taking 1 and has exactly these losing piles, with the list continuing in the same way forever. If no rule works, explain why.

\vspace{6pt}
\begin{tikzpicture}
  \foreach \t [count=\r from 0] in {{4, 8, 12, 16, 20, 24, \dots},{2, 5, 7, 10, 12, 15, \dots},{3, 7, 10, 14, 17, 21, \dots},{3, 5, 8, 10, 13, 15, \dots}}{%
    \draw[box] (0,{-\r*4.85}) rectangle ++(17.6,-4.55);
    \node[anchor=north west, font=\large] at (0.3,{-0.3-\r*4.85}) {\t};
  }
\end{tikzpicture}

\newpage
\problem Now there are two piles. Each turn takes 1 or 2 counters from one of the piles, and whoever takes the very last counter wins. In the grid, write W or L for each pair of piles with 0 to 7 counters each. Describe all the losing pairs, and explain why they are losing.

\vspace{10pt}
\hspace*{0.6cm}\pairgrid{7}{0.82}\hfill
\begin{tikzpicture}[baseline=(current bounding box.north)]\draw[box] (0,0.45) rectangle (7.6,-6.6);\end{tikzpicture}

\vspace{16pt}
\problem Two piles again, but now each turn takes 1, 3, or 4 counters from one of the piles. Write W or L for each pair of piles with 0 to 8 counters each.

\vspace{10pt}
\hspace*{0.6cm}\pairgrid{8}{0.8}

\newpage
\problem For the rule 1, 3, or 4, give each single pile a number as follows. An empty pile gets 0. Any other pile gets the smallest of the numbers 0, 1, 2, 3, \dots\ that is not the number of a pile you can move to. Find the numbers of the piles from 0 to 20. Then decide whether two piles of 11 and 13 counters are a winning pair or a losing pair, do the same for two piles of 12 and 18 counters, and explain your answers.

\vspace{4pt}
\numtable{0}{20}{11}{1.6}{0.7}{1.0}{0.35}

\vspace{6pt}
\workbox{4.4}

\vspace{14pt}
\problem Each turn takes 1, 6, or 9 counters. Write W or L under each pile from 1 to 40, and find where the pattern starts to repeat. Then explain why, for every rule, the pattern of W and L must repeat from some point on.

\vspace{4pt}
\numtable{1}{40}{20}{0.88}{0.65}{0.95}{0.35}

\vspace{6pt}
\workbox{4.0}

\end{document}
